%%%%%%%%%%%%%%%%%%%%%%%%%%%%%%%%%%%%%%%%%%%%%%%%%%%%%%%%%%%%%%%%%%%%%%%%%%%%%%%%
%%%%%%%%%%%%%%%%%%   Vorlage für eine Abschlussarbeit   %%%%%%%%%%%%%%%%%%%%%%%%
%%%%%%%%%%%%%%%%%%%%%%%%%%%%%%%%%%%%%%%%%%%%%%%%%%%%%%%%%%%%%%%%%%%%%%%%%%%%%%%%

% Erstellt von Maximilian Nöthe, <maximilian.noethe@tu-dortmund.de>
% ausgelegt für lualatex und Biblatex mit biber

% Kompilieren mit 
% latexmk --lualatex --output-directory=build thesis.tex
% oder einfach mit:
% make

\documentclass[
  tucolor,       % remove for less green,
  BCOR=12mm,     % 12mm binding corrections, adjust to fit your binding
  parskip=half,  % new paragraphs start with half line vertical space
  open=any,      % chapters start on both odd and even pages
  cleardoublepage=plain,  % no header/footer on blank pages
]{tudothesis}


% Warning, if another latex run is needed
\usepackage[aux]{rerunfilecheck}

% just list chapters and sections in the toc, not subsections or smaller
\setcounter{tocdepth}{1}

%------------------------------------------------------------------------------
%------------------------------ Fonts, Unicode, Language ----------------------
%------------------------------------------------------------------------------
\usepackage{fontspec}
\defaultfontfeatures{Ligatures=TeX}  % -- becomes en-dash etc.

% german language
\usepackage{polyglossia}
\setdefaultlanguage{english}

% for english abstract and english titles in the toc
\setotherlanguages{german}

% intelligent quotation marks, language and nesting sensitive
\usepackage[autostyle]{csquotes}

% microtypographical features, makes the text look nicer on the small scale
\usepackage{microtype}

%------------------------------------------------------------------------------
%------------------------ Math Packages and settings --------------------------
%------------------------------------------------------------------------------

\usepackage{amsmath}
\usepackage{amssymb}
\usepackage{mathtools}

% Enable Unicode-Math and follow the ISO-Standards for typesetting math
\usepackage[
  math-style=ISO,
  bold-style=ISO,
  sans-style=italic,
  nabla=upright,
  partial=upright,
]{unicode-math}
\setmathfont{Latin Modern Math}

% nice, small fracs for the text with \sfrac{}{}
\usepackage{xfrac}  


%------------------------------------------------------------------------------
%---------------------------- Numbers and Units -------------------------------
%------------------------------------------------------------------------------

\usepackage[
  locale=DE,
  separate-uncertainty=true,
  per-mode=symbol-or-fraction,
]{siunitx}
\sisetup{math-micro=\text{µ},text-micro=µ}

%------------------------------------------------------------------------------
%-------------------------------- tables  -------------------------------------
%------------------------------------------------------------------------------

\usepackage{booktabs}       % \toprule, \midrule, \bottomrule, etc

%------------------------------------------------------------------------------
%-------------------------------- graphics -------------------------------------
%------------------------------------------------------------------------------

\usepackage{graphicx}
% currently broken
% \usepackage{grffile}

% allow figures to be placed in the running text by default:
\usepackage{scrhack}
\usepackage{float}
\floatplacement{figure}{htbp}
\floatplacement{table}{htbp}

% keep figures and tables in the section
\usepackage[section, below]{placeins}


%------------------------------------------------------------------------------
%---------------------- customize list environments ---------------------------
%------------------------------------------------------------------------------

\usepackage{enumitem}

%------------------------------------------------------------------------------
%------------------------------ Bibliographie ---------------------------------
%------------------------------------------------------------------------------

\usepackage[
  backend=biber,   % use modern biber backend
  autolang=hyphen, % load hyphenation rules for if language of bibentry is not
                   % german, has to be loaded with \setotherlanguages
                   % in the references.bib use langid={en} for english sources
]{biblatex}
\addbibresource{references/simulation.bib}  % the bib file to use
\DefineBibliographyStrings{english}{andothers = {{et\,al\adddot}}}  % replace u.a. with et al.


% Last packages, do not change order or insert new packages after these ones
\usepackage[pdfusetitle, unicode, linkbordercolor=tugreen]{hyperref}
\usepackage{bookmark}
\usepackage[shortcuts]{extdash}

%------------------------------------------------------------------------------
%-------------------------    Angaben zur Arbeit   ----------------------------
%------------------------------------------------------------------------------

\author{Jan Benjamin Soedingrekso}
\title{Muons? - Muons!}
\date{2020}
\birthplace{Wuppertal}
\chair{Lehrstuhl für Experimentelle Physik V}
\division{Fakultät Physik}
\thesisclass{eines Doktors der Naturwissenschaften (Dr. rer. nat.)}
\submissiondate{31. September 2020}
\firstcorrector{Prof.~Dr.~Dr.~Wolfgang Rhode}
\secondcorrector{Prof.~Dr.~Zweitgutachter}

% tu logo on top of the titlepage
\titlehead{\includegraphics[height=1.5cm]{logos/tu-logo.pdf}}

\begin{document}
\frontmatter
\maketitle

% Gutachterseite
\makecorrectorpage

% hier beginnt der Vorspann, nummeriert in römischen Zahlen
\thispagestyle{plain}

\section*{Abstract}

Muons are the dominant particle type measured in almost every underground experiment mainly driven by the high production rate of muons in cosmic-ray induced air showers as well as the long muon range.
Due to their stochastic propagation behavior, they can remain undetected with minimal energy losses in veto regions while producing a signal-like signature with a large stochastic energy loss inside a detector.
Therefore, accurate description of theoretical models and precise treatment in simulations as well as a validation of the cross-section with measurements are required.

In this thesis, systematic uncertainties in simulations of high-energy muons were analyzed and improved, which can be divided into three parts.
The theoretical models of the cross-sections were revised and radiative corrections for the pair production interaction were calculated.
In a next step, the Monte-Carlo simulation library PROPOSAL was completely restructured in a modular design to include more accurate models and corrections.
Due to its improved usability through the modular design and its accessibility as free open-source software, PROPOSAL is now used in many applications, from large simulation frameworks, such as the CORSIKA air shower simulation, to small simulation studies.
The third part consisted of a feasibility study using PROPOSAL to measure the bremsstrahlung cross-section from the energy loss distribution, which can be measured in cubic kilometer-sized detectors.
For a detector resolution similar to that of the IceCube neutrino telescope, the bremsstrahlung normalization was estimated with an uncertainty of $\pm \SI{4}{\%}$.


\section*{Kurzfassung}
\begin{german}
Myonen sind der dominierende Teilchentyp, der in fast allen Untergrundexperimenten gemessen wird, hauptsächlich bedingt durch die hohe Produktionsrate von Myonen in durch kosmische Strahlung induzierten Luftschauern sowie die große Myonenreichweite.
Aufgrund ihres stochastischen Propagationsverhaltens können sie mit minimalen Energieverlusten unentdeckt durch Vetoregionen propagieren und innerhalb des Detektors mit einem großen stochastischen Energieverlust eine signalartige Signatur erzeugen.
Daher sind eine genaue Beschreibung der theoretischen Modelle und eine präzise Behandlung in Simulationen sowie eine Validierung des Wirkungsquerschnitts mit Messungen erforderlich.

In dieser Arbeit wurden systematische Unsicherheiten in Simulationen hochenergetischer Myonen analysiert und verbessert, was in drei Teile unterteilt werden kann.
Die theoretischen Modelle der Wirkungsquerschnitte wurden überarbeitet und Strahlungskorrekturen für die Paarproduktionswechselwirkung wurden berechnet.
In einem nächsten Schritt wurde die Monte-Carlo-Simulationsbibliothek PROPOSAL in einem modularen Design komplett umstrukturiert, um genauere Modelle und Korrekturen einbeziehen zu können.
Aufgrund der verbesserten Nutzbarkeit durch den modularen Aufbau und der Zugänglichkeit als freie Open-Source-Software wird PROPOSAL inzwischen in vielen Anwendungen eingesetzt, von großen Simulations-Frameworks, wie der Luftschauer-Simulation CORSIKA, bis hin zu kleinen Simulationsstudien.
Der dritte Teil bestand aus einer Machbarkeitsstudie unter Verwendung von PROPOSAL zur Messung des Bremsstrahlungsquerschnitts aus der Energieverlustverteilung, die in kubikkilometergroßen Detektoren gemessen werden kann.
Für eine Detektorauflösung, welcher der des IceCube-Neutrinoteleskops ähnelt, wurde die Bremsstrahlungsnormalisierung mit einer Unsicherheit von $\pm \SI{4}{\%}$ abgeschätzt.

\end{german}

\tableofcontents

\mainmatter
% Hier beginnt der Inhalt mit Seite 1 in arabischen Ziffern
\chapter{Introduction}

Muons have been first discovered in cloud chamber observations in 1936 \cite{Anderson36}.
Due to their propagated range, energy loss and deflection their signature in the detector didn't match the behaviour of an electron or a proton.
In particular it's the range muons can proapagte through large volumes of media, that makes them special and interesting for nearly all particle detectors on earth.
Muons are the only particle type from cosmic ray induced air showers that reach the surface of the earth at sea leavel.
With the high rate of cosmic rays hitting the atmosphere, their secondary muons contribute to a third of the natural radiation consumption on earth.
Even for underground experiments muons are the dominant measured event signature.

Regarding muon energies above a GeV, most muons detected at the earths surface coming from the sky and are secondary products of cosmic rays interactions.
For muons coming from down the earth, propagating horizontical or upwards, originate from neutrino interactions.
As neutrinos can propagate through the earth they can interact before the detector or inside and produce their leptonic partner, e.g. muons.

For Cosmic Ray Air shower detectors and neutrino detectors it is important to have a good understanding of the muon physics.
Next to the cross sections, also the remaining uncertainties and an understanding of their effects in analysis is important.

Since, there is no astrophysical source, sending a test beam of cosmic rays or neutrinos for all energies to calibrate the detection, Monte Carlo simulations are necessary to understand the measured data.
% Regarding the electroweak interactions of leptons and neutrinos, PROPOSAL is a simulation library for Monte-Carlo simulation providing multiple interaction and scattering models

In Chapter \ref{sec:theory} the theory is presented.
\chapter{Muons in Astroparticlephysics}
\label{sec:theory}

Slide 5 and more

\section{Muons from Cosmic Rays}

Air shower

\section{Muons from Neutrinos}

IceCube

\section{Muon Simulation}

PROPOSAL \cite{JHK, PROPOSAL, dPROPOSAL}

\appendix
% Hier beginnt der Anhang, nummeriert in lateinischen Buchstaben
\chapter{Appendix}


\section{Tables of the Photonuclear Interactions} \label{sec:photo_tables}

\begin{table}
    \centering
    \caption{Measured Photon-Nucleon cross-section $\sigma_{\gamma N}$ used in the photonuclear cross-section of the Rhode parametrization\cite{Rhode93PhD}, described in \secref{sec:photo_vmd}. The the photon energies $E_\gamma$ are equally distributed in $\log_{10}$ space.}
    \label{tab:photo_rhode_data}
    \sisetup{table-format = 1.4e1}
    \begin{tabular}{S[table-format=1.4e+1] S[table-format=3.6] | S S[table-format=3.2] | S S[table-format=3.2]}
        \toprule
        {$E_\gamma$/GeV} & {$\sigma_{\gamma N}/\mu$b} & {$E_\gamma$/GeV} & {$\sigma_{\gamma N}/\mu$b} & {$E_\gamma$/GeV} & {$\sigma_{\gamma N}/\mu$b} \\
        \midrule
        1.0000e-01 & 0.066667 & 5.2481e+01 & 114.37 & 2.7542e+04 & 211.78 \\
        1.4454e-01 & 0.096363 & 7.5858e+01 & 114.79 & 3.9811e+04 & 223.50 \\
        2.0893e-01 & 159.74 & 1.0965e+02 & 115.86 & 5.7544e+04 & 235.88 \\
        3.0200e-01 & 508.10 & 1.5849e+02 & 117.61 & 8.3176e+04 & 248.92 \\
        4.3652e-01 & 215.77 & 2.2909e+02 & 120.01 & 1.2023e+05 & 262.63 \\
        6.3096e-01 & 236.40 & 3.3113e+02 & 123.08 & 1.7378e+05 & 277.01 \\
        9.1201e-01 & 201.92 & 4.7863e+02 & 126.81 & 2.5119e+05 & 292.05 \\
        1.3183e+00 & 151.38 & 6.9183e+02 & 131.21 & 3.6308e+05 & 307.75 \\
        1.9055e+00 & 145.41 & 1.0000e+03 & 136.28 & 5.2481e+05 & 324.12 \\
        2.7542e+00 & 132.10 & 1.4454e+03 & 142.01 & 7.5858e+05 & 341.16 \\
        3.9811e+00 & 128.55 & 2.0893e+03 & 148.40 & 1.0965e+06 & 358.86 \\
        5.7544e+00 & 125.05 & 3.0200e+03 & 155.46 & 1.5849e+06 & 377.22 \\
        8.3176e+00 & 121.86 & 4.3652e+03 & 163.19 & 2.2909e+06 & 396.25 \\
        1.2023e+01 & 119.16 & 6.3096e+03 & 171.57 & 3.3113e+06 & 415.95 \\
        1.7378e+01 & 117.02 & 9.1201e+03 & 180.63 & 4.7863e+06 & 436.31 \\
        2.5119e+01 & 115.50 & 1.3183e+04 & 190.35 & 6.9183e+06 & 457.33 \\
        3.6308e+01 & 114.61 & 1.9055e+04 & 200.73 & 1.0000e+07 & 479.02 \\
        \bottomrule
    \end{tabular}
\end{table}

\begin{table}
    \caption{Fit values of the ALLM parametrization estimated in ALLM91\cite{Abramowicz91}, ALLM97 \cite{Abramowicz97}, and AbtFT\cite{Abt17PhotoQ2}, described in \secref{sec:photo_regge}.}
    \label{tab:photo_q2_params}
    \begin{center}
    \begin{tabular}{l | S[table-format=+2.5] S[table-format=+2.5] S[table-format=+2.4]}
        \toprule
        Parameters & {ALLM91} & {ALLM97} & {AbtFT17} \\
        \midrule
        $a_{1,R}$ & 0.60408 & 0.58400 & 0.882 \\
        $a_{2,R}$ & 0.17353 & 0.37888 & 0.082 \\
        $a_{3,R}$ & 1.61812 & 2.6063 & -8.5 \\
        $b_{1,R}$ & 1.26066 & 0.01147 & 0.339 \\
        $b_{2,R}$ & 1.83624 & 3.7582 & 3.38 \\
        $b_{3,R}$ & 0.81141 & 0.49338 & 1.07 \\
        $c_{1,R}$ & 0.67639 & 0.80107 & -0.636 \\
        $c_{2,R}$ & 0.49027 & 0.97307 & 3.37 \\
        $c_{3,R}$ & 2.66275 & 3.4942 & -0.660 \\
        \midrule
        $a_{1,P}$ & -0.04503 & -0.0808 & -0.075 \\
        $a_{2,P}$ & -0.36407 & -0.44812 & -0.470 \\
        $a_{3,P}$ & 8.17091 & 1.1709 & 9.2 \\
        $b_{1,P}$ & 0.49222 & 0.36292 & -0.477 \\
        $b_{2,P}$ & 0.52116 & 1.8917 & 54.0 \\
        $b_{3,P}$ & 3.55115 & 1.8439 & 0.073 \\
        $c_{1,P}$ & 0.26550 & 0.28067 & 0.356 \\
        $c_{2,P}$ & 0.04856 & 0.22291 & 0.171 \\
        $c_{3,P}$ & 1.04682 & 2.1979 & 18.6 \\
        \midrule
        $m_\gamma^2$ / \si{GeV^2} & 0.30508 & 0.31985 & 0.388 \\
        $m_R^2$ / \si{GeV^2} & 0.20623 & 0.15052 & 0.838 \\
        $m_P^2$ / \si{GeV^2} & 10.67564 & 49.457 & 50.8 \\
        $\Lambda^2$ / \si{GeV^2} & 0.06527 & 0.06527 & {\num{4.4e-9}} \\
        $(Q_0^2 - \Lambda^2)$ / \si{GeV^2} & 0.27799 & 0.52544 & {\num{1.87e-5}} \\
        \bottomrule
    \end{tabular}
    \end{center}
\end{table}

%%%
%%%
%%%

\section{Radiative Corrections to the Pair Production} \label{sec:epair_nlo_calc}

In this section, radiative corrections of the pair production for the muon line are calculated, as described in \secref{sec:epair_rad_corr}, considering only the $e$-diagrams, since they are the dominating process.

Before describing the calculation, the used dilogarithm is defined by
\begin{align} \label{eq:dilog}
    \dilog(x) = - \mathrm{Re} \int_{0}^{x} \frac{\logn(1 - t)}{t} \dif t .
\end{align}

The tree-level cross-section of pair production and bremsstrahlung calculated by Bugaev \cite{Bugaev77} are used as base to estimate radiative corrections.
Due to the modular structure of the parametrization, dividing the cross-section into the muon line, the electron line, and the nuclear Interaction as three distinct tensors, each tensor can easily replaced introducing e.g. a correction to a fermion line.
Furthermore, the bremsstrahlung cross-section is calculated in the same formalism allowing a combination of both processes due the their consistent treatment.

The differential cross section is given by
\begin{align}
    \frac{\dif \sigma}{\dif \omega \dif \varepsilon_+ \dif k^2 \dif q^2 \dif \nu} =
        \frac{Z^2 \alpha^4 m_e^2}{16 \pi^2 p^2} \frac{1}{k^2 q^2}
        (W_1(L_1 M_{\alpha \alpha}^{\mu \mu} - L_2 M_{44}{\mu \mu})
        + W_2 (L_2 M_{44}^{44} - L_1 M_{\alpha \alpha}^{44})) ,
\end{align}
where $k^2$ and $q^2$ are the squared momentum and $\omega and \nu$ the energies of the virtual photons connecting the electron with the muon line and the electron line with the nucleus respectively.
The tensors of the nucleus $W$, the electron line $M$ and the muon line $L$ are already contracted.
Instead of being differential in the asymmetry $\rho$, this cross section uses the energy of the positron $\varepsilon_+$.

To simplify the calculation, the same relativistic approximations as described in \cite{Bugaev77} are applied, consisting of
\begin{itemize}
    \item a high $\gamma$ factor of the muon and the electron-positron pair ($E_\mu^2 \gg m_\mu^2$ and $\varepsilon_+^2 \gg m_e^2$)
    \item small momentum transfers ($k^2, q^2 \ll E_\mu^2$)
    \item not the smallest energy losses ($\omega^2 \gg m_\mu^2$)
\end{itemize}
Furthermore, the nucleus is assumed to be heavy and have no recoil, thus neglecting the energy transfered to the nucleus $\nu \approx 0$.
Therefore, also the spatial components of the electromagnetic structure functions are neglected $W_1 \approx 0$, while $W_2$ is defined by
\begin{align}
    W_2 (q^2) = -((f_n - f_a)^2 + 1/Z*(f_{n,\mathrm{in}} + f_{a,\mathrm{in}})) .
\end{align}
The form factors are taken from \cite{Tsai74, Tsai77, Andreev94Brems}
\begin{subequations}
\begin{align}
    f_n   &= (1 + a^2 q^2/12)^{-2}, % alternatively = Fn_fermi(q2)
    & f_a &= (1 + b^2 q^2)^{-1}, \\
    f_{n,\mathrm{in}}   &= 1 - f_n^2,
    & f_{a,\mathrm{in}} &= \left(\frac{c^2 q^2}{1 + c^2 q^2} \right)^2 ,
\end{align}
\end{subequations}
with
\begin{subequations}
\begin{align}
    a &= (0.58 + 0.82 A^{\sfrac13}) \cdot 5.07/\si{GeV},
    \\
    b &= \frac{B_{\mathrm{el}}}{\sqrt{e} m_e} Z^{-\sfrac13},
    \\
    c &= \frac{B_{\mathrm{inel}}}{\sqrt{e} m_e} Z^{-\sfrac23}.
\end{align}
\end{subequations}

For the tree-level process, the functions for the muon line $L_{1,2}$ can be described with eq. (44) of \cite{Bugaev77}
% \begin{align}
%     L_1(k^2) = \frac{2}{k^{\ast}} \left( \frac{k^2}{2} - m_\mu^2 - \frac{k^2}{{k^{\ast}}^2} ( E_\mu E_{\mu} + k^2 / 4) \right)
%     \qquad
%     L_2(k^2) = \frac{2 k^2}{k^{\ast}^3} \left( \frac{k^2}{2} - m_\mu^2 - 3 \frac{k^2}{{k^{\ast}}^2} ( E_\mu E_{\mu} + k^2 / 4) \right)
% \end{align}
and tensors of the electron line $M_{44}^{44}$ and $M_{\alpha \alpha}^{44}$ with eq. (36) of \cite{Bugaev77}.
Integrating the tree-level diagram over $q^2, k^2$, and $\varepsilon_+$ results in a similar differential cross-section in the energy loss, compared to the $e$-diagram of Kelner et al. (c.f. \secref{sec:epair_screen_approx}), as shown in \figref{fig:epair_all_compare}.
\begin{figure}
    \centering
    \includegraphics[width=0.8\textwidth]{./plots/radiative_corrections/plot_epair_dsigma_loop3.pdf}
    \caption{Differential cross section of the pair production cross section including the tree-level of the $e$-diagram (Bugaev) and its radiative corrections along the muon line, vacuum polarization, vertex correction, and soft and hard bremsstrahlung. The $e$-diagram of the Kelner et al. parametrization is included and compared to the calculated $e$.
    Also the $\mu$-diagram of Kelner et al. is included.}
    \label{fig:epair_all_compare}
\end{figure}

The overall correction due to the radiative corrections are in the percent level. However, regarding the differential cross-section, this correction can have larger effect, especially at higher energy losses.
This is mainly driven by the hard bremsstrahlung with an energy loss spectrum similar to the $\mu$-diagram.
In future works, this only numerically integrated contribution needs to parameterized to be usable in simulations like PROPOSAL, as the correction is on the order as the $\mu$-diagram.

In the following subsections, the calculation of the radiative correction is described.

\subsection{Vacuum Polarization}

Out of the tree-level cross-section, the vacuum polarization can simply be included with a correction factor before the integration over $k^2$, since it only depends on $k^2$.
The correction factor to the tree process was calculated in \cite{Mork65} eq. (IV.2/3) and without setting the electron mass to 1 this results in
\begin{align}
s_{\mathrm{vac}}(k^2) = -\frac{2 \alpha}{3 \pi}
        \left[ \frac53 - \zeta - \left( 1-\frac{\zeta}{2} \right) \sqrt{1+\zeta} 
        \logn \left( \frac{k^2 + k^2 \sqrt{1 + \zeta}}{2m_e^2} + 1 \right) \right]
\end{align}
using
\begin{align}
    \zeta = \frac{4 m_e^2}{k^2} .
\end{align}

\subsubsection{Vertex Correction and Soft Bremsstrahlung}

The Vertex correction to a fermion line for an outgoing space-like photon has been calculated in \cite{Akhiezer81} in eq. (5.1.37) leading to the renormalized form
\begin{align}
    \Lambda_\mu = \frac{\alpha}{\pi} \left( \gamma_\mu \Lambda_1 + i \frac{\gamma_\mu \hat{k} - \hat{k} \gamma_\mu}{8 M} \Lambda_2 \right) ,
\end{align}
with the expressions
\begin{subequations} \label{eq:lambda_vertex}
\begin{align}
    \Lambda_1 &= \left( \logn \frac{M}{\lambda} - 1 \right)
        \left( \frac{1 + a^2}{2a} \logn b + 1 \right)
        + \frac{\logn b}{4a} ,
        \\
        &\quad - \frac{1 + a^2}{4a} \left( -\frac{\pi^2}{6} + \frac{\logn^2 b}{2} - 2 \logn b \logn(1 + b) - 2 \dilog (-b) \right) \nonumber \\
    \Lambda_2 &= \frac{a^2 - 1}{2a} \logn b ,
\end{align}
\end{subequations}
using abbreviations
\begin{align}
    a = \sqrt{1 + \frac{4 M^2}{k^2}}, \qquad
    b = \frac{a - 1}{a + 1}.
\end{align}
The tensor of the muon line including the vertex correction can be written as 
\begin{align}
    L_{\mu\nu} = \Lambda_{\mu} (i \hat{p}_1 - M) \gamma_\nu (i \hat{p}_2 - M) .
\end{align}
To contract these tensors, and calculating the traces over the gamma matrices, the computer algebra program FORM \cite{Vermaseren00Form} was used resulting in
\begin{subequations} \label{eq:vacuum_l12}
\begin{align}
    \delta^{\mu\nu} L_{\mu\nu} = L_{\mu\mu} &= \left( \frac{k^2}{2} - M^2 \right) \Lambda_1 - 3 k^2 \Lambda_2, \\
    \frac{1}{M^2} p^\mu p^\nu L_{\mu\nu} = L_{44} &= \left( \frac{k^2}{4} - E_1 E_2 \right)\Lambda_1 - {k^{\ast}}^2 \Lambda_2 .
\end{align}
\end{subequations}
Including \eqref{eq:vacuum_l12} into the definition of $L_{1,2}$ in eq. (9) of \cite{Bugaev77}, eq. (44) of \cite{Bugaev77} changes to
\begin{subequations}
\begin{align}
    L_1 &= \frac{2}{k^\ast} L_{\mu\mu} - \frac{k^2}{{k^{\ast}}^2} L_{44}, \\
    L_2 &= \frac{2k^2}{{k^{\ast}}^3} L_{\mu\nu} - 3\frac{k^2}{{k^{\ast}}^2} L_{44} .
\end{align}
\end{subequations}
In \eqref{eq:lambda_vertex} an effective photon mass $\lambda$ was introduced to avoid the infrared divergence of this NLO process.

This unphysical photon mass get canceled out by including also a soft bremsstrahlung contribution using the same photon mass.
For the soft bremsstrahlung, the correction calculated in \cite{Mork65} is used
\begin{equation}
\begin{aligned}
    I_{\text{soft}} &= 2(1 - 2y \coth (2y) \logn \frac{E_{\gamma, \text{max soft}} m_\mu}{\lambda \sqrt{E_1 E_2}}
    \\
    &\quad + \frac12 \coth (2y) \left[
        2y \logn \frac{(1 - \zeta^2)(\zeta^2\eta^2 - 1)}{4 \zeta^2 (\eta^2 - 1)}
        + \logn \frac{\eta + 1}{\eta - 1} \logn \frac{\zeta \eta + 1}{\zeta \eta - 1} \right.
    \\
    &\quad \hphantom{+ \frac12 \coth (2y) +}
        + \dilog \left( \frac{(1 + \zeta)(\zeta \eta + 1)}{2 \zeta (\eta + 1)} \right)
        - \dilog \left( \frac{(1 - \zeta)(\zeta \eta - 1)}{2 \zeta (\eta + 1)} \right)
    \\
    &\quad \hphantom{+ \frac12 \coth (2y) +} \left.
        + \dilog \left( \frac{(1 + \zeta)(\zeta \eta - 1)}{2 \zeta (\eta - 1)} \right)
        - \dilog \left( \frac{(1 - \zeta)(\zeta \eta - 1)}{2 \zeta (\eta - 1)} \right)
    \right] ,
\end{aligned}
\end{equation}
with
\begin{align}
    y = \operatorname{arsinh} \sqrt{\frac14 \frac{k^2}{M^2}} ,
    \qquad
    \zeta = \tanh y ,
    \qquad
    \eta = \frac{E_1 + E_2}{-E_1 + E_2}
\end{align}
since it only valid for small photon energies, a maximum photon energy $E_{\gamma, \text{max soft}}$ is introduced.

\subsection{Hard Bremsstrahlung}

The bremsstrahlung cross-section can be derived from the $\mu$-diagram by setting $k^2 = 0$.
To include the hard bremsstrahlung into the $e$-diagram, the muon line of the $e$-diagram $L_{1,2}$ gets replace by the muon line of the bremsstrahlung $M_{\alpha \alpha}^{44}$ and $M_{\alpha \alpha}^{\mu \mu}$.
Compared to described in eq. (49) of \cite{Bugaev77}, the substitutions $\omega \to \omega_{\mathrm{brems}}$, $\nu \to -\omega$, and $q^2 \to k^2$ is necessary.
Now, the differential cross section is additionally differential in the energy of the hard bremsstrahlung photon $\omega_{\mathrm{brems}}$ with the lower integration limit defined by the maximum energy of the soft photon $E_{\gamma, \text{max soft}}$.



\backmatter
\printbibliography

\cleardoublepage
\thispagestyle{empty}
\section*{Eidesstattliche Versicherung}
Ich versichere hiermit an Eides statt, dass ich die vorliegende Abschlussarbeit mit dem Titel \enquote{\thetitle} selbstständig und ohne unzulässige fremde Hilfe erbracht habe.
Ich habe keine anderen als die angegebenen Quellen und Hilfsmittel benutzt, sowie wörtliche und sinngemäße Zitate kenntlich gemacht. 
Die Arbeit hat in gleicher oder ähnlicher Form noch keiner Prüfungsbehörde vorgelegen.

\vspace*{1cm}\noindent
\begin{center}
  \begin{tabular}{@{}p{0.4\textwidth}@{\hspace{0.15\textwidth}}p{0.4\textwidth}@{}}
  \rule{\linewidth}{0.25pt}& \rule{\linewidth}{0.25pt}\\
  Ort, Datum & Unterschrift
  \end{tabular}
\end{center}

\subsection*{Belehrung}
Wer vorsätzlich gegen eine die Täuschung über Prüfungsleistungen betreffende Regelung einer Hochschulprüfungsordnung verstößt, handelt ordnungswidrig.
Die Ordnungswidrigkeit kann mit einer Geldbuße von bis zu \SI[round-mode=places, round-precision=2]{50000}{€} geahndet werden. 
Zuständige Verwaltungsbehörde für die Verfolgung und Ahndung von Ordnungswidrigkeiten ist der Kanzler/die Kanzlerin der Technischen Universität Dortmund. 
Im Falle eines mehrfachen oder sonstigen schwerwiegenden Täuschungsversuches kann der Prüfling zudem exmatrikuliert werden \mbox{(\S\,63 Abs. 5 Hochschulgesetz --HG--).}

Die Abgabe einer falschen Versicherung an Eides statt wird mit Freiheitsstrafe bis zu 3 Jahren oder mit Geldstrafe bestraft.

Die Technische Universität Dortmund wird ggf.\ elektronische Vergleichswerkzeuge (wie z.\,B.\ die Software \enquote{turnitin}) zur Überprüfung von Ordnungswidrigkeiten in Prüfungsverfahren nutzen. \\[\baselineskip]

\noindent Die oben stehende Belehrung habe ich zur Kenntnis genommen.\\[1cm]
\begin{center}
\begin{tabular}{@{}p{0.4\textwidth}@{\hspace{0.15\textwidth}}p{0.4\textwidth}@{}}
\rule{\linewidth}{0.25pt}& \rule{\linewidth}{0.25pt}\\
Ort, Datum & Unterschrift
\end{tabular}
\end{center}

\end{document}
